\documentclass[12pt]{article}
\usepackage[margin=1in]{geometry}
\usepackage{amsmath, amssymb, amsfonts}
\usepackage{enumitem}
\usepackage{booktabs}
\usepackage{hyperref}
\hypersetup{colorlinks=true, linkcolor=blue, urlcolor=blue}

\setlist[enumerate]{leftmargin=*}

\title{\textbf{Solutions -- Quiz 1} \vspace{ 0.4cm }\\ECON 101: Macroeconomics }
\author{Instructor: Muhebullah Karimzada}
\date{January 22,  2026}

\begin{document}
\maketitle

\section*{Instructions}
Answer all questions. Show all work for numerical problems. Multiple-choice questions have one correct answer.

\section*{Part 1: Multiple Choice (2 points each)}

\begin{enumerate}
    \item Economics is best defined as the study of:
    \begin{enumerate}[label=\Alph*.]
        \item How governments set tax rates  
        \item How to eliminate scarcity  
        \item The choices people make to attain their goals, given scarce resources  
        \item How firms maximize profits only  
    \end{enumerate}

    \item Which of the following is an example of a \textbf{positive economic statement}?  
    \begin{enumerate}[label=\Alph*.]
        \item The Bank of Canada should lower interest rates.  
        \item A higher minimum wage will lead to higher unemployment among teenagers.  
        \item Canada ought to spend more on public health care.  
        \item Equity should be prioritized over efficiency.  
    \end{enumerate}

    \item In the economic problem of scarcity:  
    \begin{enumerate}[label=\Alph*.]
        \item Wants are limited, resources are unlimited.  
        \item Both wants and resources are unlimited.  
        \item Wants are unlimited, resources are limited.  
        \item Both wants and resources are limited.  
    \end{enumerate}
\newpage
    \item The \textbf{Production Possibilities Frontier (PPF)} is bowed outward because:
    \begin{enumerate}[label=\Alph*.]
        \item Opportunity costs are constant.  
        \item Resources are not equally productive in all uses.  
        \item Scarcity does not apply.  
        \item Technology is fixed.  
    \end{enumerate}

    \item If Canada can produce wheat more efficiently than Japan, while Japan can produce electronics more efficiently than Canada, trade between the two countries can be mutually beneficial due to:
    \begin{enumerate}[label=\Alph*.]
        \item Absolute advantage  
        \item Comparative advantage  
        \item Specialization without opportunity costs  
        \item Increasing marginal costs  
    \end{enumerate}
\end{enumerate}



\section*{Part 2: Numerical Questions}

\begin{enumerate}[resume]
    \item Suppose a student has 10 hours per week to allocate between working (earning \$15/hour) and studying. What is the \textbf{opportunity cost} of one additional hour of studying? \textbf{(2 points)}

    \item Consider the following simplified PPF data:

    \begin{center}
    \begin{tabular}{ccc}
    \toprule
    Choice & Cars & Computers \\
    \midrule
    A & 0 & 100 \\
    B & 20 & 80 \\
    C & 40 & 60 \\
    D & 60 & 40 \\
    E & 80 & 20 \\
    F & 100 & 0 \\
    \bottomrule
    \end{tabular}
    \end{center}

    \begin{enumerate}
        \item What is the opportunity cost of increasing car production from 20 to 40?  \textbf{(2 points)}
        \item Does this country face increasing, decreasing, or constant opportunity costs? Explain.\textbf{(2 points)}
    \end{enumerate}

    \item Country A can produce 50 tons of steel or 100 tons of corn with all its resources. Country B can produce 60 tons of steel or 120 tons of corn.
    \begin{enumerate}
        \item Which country has the \textbf{absolute advantage} in each good?  \textbf{(2 points)}
        \item Which country has the \textbf{comparative advantage} in steel?  \textbf{(2 points)}
    \end{enumerate}
\end{enumerate}

\newpage
\section*{Answer Key}

\textbf{Part 1: Multiple Choice} \\
1. C \\
2. B \\
3. C \\
4. B \\
5. B \\


\noindent \textbf{Part 2: Numerical} \\
6. Opportunity cost of 1 hour studying = \$15 lost wages. \\
7a. Going from 20→40 cars reduces computers from 80→60 → opp. cost = 20 computers. \\
7b. Opportunity cost is constant (linear trade-off). \\
8a. Absolute advantage: steel = B, corn = B. \\
8b. Comparative advantage: compute opp. cost. A: 1 steel = 2 corn. B: 1 steel = 2 corn. Both same $\Rightarrow$ no comparative advantage (special case). \\

\end{document}
